\begin{table}[H]
\centering
\caption{Κύρια χαρακτηριστικά \ENG{Sequencing} \& \ENG{Navigation} στο \ENG{SCORM} 2004}
\label{tab:scorm2004_sn_features}
\small
\setlength{\tabcolsep}{5pt}
\renewcommand{\arraystretch}{1.25}
\begin{chaptertablebox}
\begin{tabular}{|p{4.3cm}|p{9.2cm}|}
\hline
\textbf{Χαρακτηριστικό} & \textbf{Περιγραφή} \\
\hline
\hline
\textbf{\ENG{Activity} \ENG{Trees}} & Ιεραρχική οργάνωση εκπαιδευτικών δραστηριοτήτων σε δενδρική δομή. \\
\hline
\textbf{\ENG{Sequencing} \ENG{Rules}} & Κανόνες προϋποθέσεων και μετα-συνθηκών για πρόσβαση σε δραστηριότητες. \\
\hline
\textbf{\ENG{Rollup} \ENG{Rules}} & Κανόνες που καθορίζουν πώς η \ENG{completion/success} των θυγατρικών επηρεάζει την κατάσταση της γονικής δραστηριότητας. \\
\hline
\textbf{\ENG{Navigation} \ENG{Requests}} & Δυνατότητα αιτημάτων πλοήγησης (π.χ. «επόμενο», «προηγούμενο», «πήγαινε σε συγκεκριμένη δραστηριότητα»). \\
\hline
\textbf{\ENG{Control} \ENG{Modes}} & Ρυθμίσεις που καθορίζουν τον τρόπο πλοήγησης (π.χ. \texttt{\ENG{choice}}, \texttt{\ENG{flow}}). \\
\hline
\end{tabular}
\end{chaptertablebox}
\end{table}
